\section{Diffusion Policy, RDT y la línea de los modelos del mundo}

Los capítulos anteriores giraron principalmente en torno a Transformer/VLM/VLA; este capítulo examina la línea de generación de action. Diffusion Policy aplica los modelos de difusión al aprendizaje de políticas visomotoras: no genera imágenes, sino trayectorias de acciones. Esto resulta natural para los robots, porque una secuencia de acciones también puede verse como un objeto que debe eliminarse de ruido y generarse paso a paso.

\begin{figure}[H]
\centering
\includegraphics[width=0.96\textwidth]{ch30-06320.jpg}
\caption{Diffusion Policy: aprendizaje de una visuomotor policy mediante difusión de acciones. Fotograma de la pantalla proyectada, aprox. 01:45:20.}
\end{figure}

\begin{knowledgebox}{Lectura de la figura: los modelos de difusión también pueden generar action}
La figura muestra los procesos de entrenamiento y de inferencia de Diffusion Policy. La transferencia central es esta: los modelos de difusión no solo sirven para generar imágenes, sino también para generar action trajectory, y son especialmente adecuados para acciones continuas y trayectorias multimodales.
\end{knowledgebox}

\subsection{RDT-1B: un modelo fundacional de difusión para manipulación bimanual}

Esta sección sigue indagando cómo debe modelarse la action misma. La línea VLM+Action Head trata el módulo de acción como un controlador de alta frecuencia, pero aún hace falta un mecanismo generativo más adecuado para trayectorias continuas, elecciones multimodales y coordinación de ambas manos. La importancia de RDT-1B consiste en llevar la intuición de Diffusion Policy a la escala de un modelo fundacional, de modo que la «generación de acciones» deje de ser solo un componente accesorio de una pequeña red de política.

RDT-1B extiende Diffusion Policy a un modelo fundacional más grande para manipulación bimanual. Busca resolver el unified action space, las distintas action head y la manipulation con dos brazos. Su importancia reside en llevar la difusión de acciones propia de los modelos pequeños hacia una escala mayor, más cercana a un foundation model.

\begin{figure}[H]
\centering
\includegraphics[width=0.96\textwidth]{ch31-06571.jpg}
\caption{RDT-1B: un modelo fundacional de difusión para robots de manipulación bimanual. Fotograma de la pantalla proyectada, aprox. 01:49:31.}
\end{figure}

\subsection{Prediction with Action y VPP}

Prediction with Action coloca la predicción visual y la predicción de acciones en un proceso conjunto de eliminación de ruido, con el fin de conectar los modelos del mundo con los VLA. Su continuación, VPP, pone el énfasis en usar predictive visual representations para formar una generalist robot policy. Lo central aquí no es solo producir la action, sino comprender a la vez los cambios visuales futuros y las consecuencias de las acciones.

\begin{figure}[H]
\centering
\includegraphics[width=0.96\textwidth]{ch32-07061.jpg}
\caption{Prediction with Action / VPP: aprendizaje de políticas visuales mediante eliminación de ruido conjunta. Fotograma de la pantalla proyectada, aprox. 01:57:41.}
\end{figure}

\begin{knowledgebox}{Lectura de la figura: los modelos del mundo entran en los VLA}
La figura incluye la predicción de video, las acciones y el proceso de eliminación de ruido. Al leerla hay que retener una idea: el robot no solo necesita saber cuál es la siguiente acción, sino también predecir cómo esa acción cambiará el mundo visual; esa es la línea de los modelos del mundo.
\end{knowledgebox}

\subsection{Direcciones futuras: UP-VLA y aprendizaje por refuerzo en línea}

Esta sección pasa de los sistemas existentes a los problemas de la siguiente etapa. Los VLA actuales ya pueden conectar visión, lenguaje y acción, pero todavía no resuelven de forma estable la unificación entre «comprender el mundo» y «predecir las consecuencias de las acciones», ni el ciclo cerrado de mejora continua en entornos reales. Por ello, Chen Jianyu (陈建宇) condensa las direcciones futuras en dos líneas: una estructura de modelo más unificada y un aprendizaje en línea más seguro y controlable.

Por último, Chen Jianyu enumera dos direcciones futuras. UP-VLA busca unificar understanding y prediction, integrando la comprensión y la predicción en un solo modelo del embodied agent. La otra dirección es el online reinforcement learning, que usa la retroalimentación en línea para seguir mejorando el VLA. La dificultad es que, si se aplica RL directamente a todo el VLM/VLA, el entrenamiento puede volverse inestable; una estrategia consiste en congelar primero el VLM, entrenar solo la action head y después liberarlo de manera gradual.

\begin{figure}[H]
\centering
\includegraphics[width=0.96\textwidth]{ch33-07592.jpg}
\caption{iRe-VLA / Online RL: mejora de un Vision-Language-Action Model mediante aprendizaje por refuerzo en línea. Fotograma de la pantalla proyectada, aprox. 02:06:32.}
\end{figure}

\begin{importantbox}{Dos líneas futuras}
La primera línea es unificar comprensión, predicción y acción, de modo que el modelo no solo entienda el mundo, sino que también pueda predecir las consecuencias de sus acciones; la segunda línea es el aprendizaje por refuerzo en línea, que permite al robot mejorar continuamente a partir de la retroalimentación de entornos reales o simulados.
\end{importantbox}

\subsection{Resumen de la sección}

Diffusion Policy y RDT llevan el modelado de la generación de acciones hacia un control continuo más potente; Prediction with Action/VPP introduce los modelos del mundo en los VLA; UP-VLA y el RL en línea apuntan hacia los modelos unificados y el aprendizaje continuo del futuro.
